\noindent
\graphicspath{{./19-07/Images/}} 
%%%%%%%%%%%%%%%%
% Page de gauche
%%%%%%%%%%%%%%%%
\begin{minipage}[t][\textheight][t]{0.33\textwidth}
    \colorbox{StrongBlue}{
        \parbox[c][\textheight][c]{\textwidth}{ 
            \color{OffWhite} 
            \begin{center}
                \begin{minipage}[t][0.6\textheight][c]{0.8\textwidth}
                    {\huge \textbf{Dimanche 19 juillet}}\\
                    \normalsize
                    \begin{center}
                        \begin{tabular}{@{} l @{\hspace{0.5cm}} l @{}}
                            Abri de départ : &\hfill L'herbaudière \\ 
                            Abri d'arrivée : &\hfill Pornic \\ 
                            Distance : &\hfill 10 Milles \\
                            Durée : &\hfill 2.5 Heures \\  
                            Départ : &\hfill  \\ 
                            Arrivée prévue : &\hfill  \\ 
                            Lever soleil 
\includegraphics[width=8px,height=8px]{../../Pictos/soleil_OffWhite.png} : &\hfill 06 h 32 \\ 
                            Coucher soleil 
\includegraphics[width=8px,height=8px]{../../Pictos/lune_OffWhite.png} : &\hfill 21 h 57 \\ 
                            \\
                            Abris replis : &  \\
                                    & • Le Palais \\
                                    & • Saint Gildas \\
                                    & • Port Haliguen \\                    
                                
                        \end{tabular}
                    \end{center}
                    
                \end{minipage}%

                \begin{minipage}[t][0.4\textheight][t]{\textwidth}
                    \begin{center}
    {\large \textit{Marées \textbf{L'Herbaudière}}}\\ \vspace{0.5em}
    
    \renewcommand{\arraystretch}{1.40} % Adjust the row height by a factor of 1.4
 \begin{tabular}{ |C{5em}|c|C{4em}|C{4em}| } \hline
    & Coef & Heure & Hauteur\\ \hline
    PM &  \multirow{2}{1em}{81} & 08 h 27 & 5.20 m\\ \cline{1-1} \cline{3-4}
    BM &                        & 15 h 19 & 1.39 m\\ \hline
    BM & \multirow{1}{1em}{75}  & 20 h 41 & 5.32 m\\ \hline
\end{tabular}
\end{center}

                    \begin{center}
    {\large \textit{Marées \textbf{Pornic}}}\\ \vspace{0.5em}
    
    \renewcommand{\arraystretch}{1.40} % Adjust the row height by a factor of 1.4
    \begin{tabular}{ |C{5em}|c|C{4em}|C{4em}| } \hline
        & Coef & Heure & Hauteur\\ \hline
    PM &  \multirow{2}{1em}{91} & 08 h 25 & 5.34 m\\ \cline{1-1} \cline{3-4}
    BM &                        & 15 h 26 & 1.34 m\\ \hline
    BM & \multirow{1}{1em}{86}  & 20 h 39 & 5.53 m\\ \hline
    \end{tabular}
\end{center}

                \end{minipage}%
            \end{center}       
        }      
    }
\end{minipage}%
\begin{minipage}[c][\textheight][t]{0.10\textwidth}
\hfill
\end{minipage}%
\begin{minipage}[c][\textheight][t]{0.50\textwidth}
    \vspace{2cm}
    \color{StrongBlue}
            {\huge \textbf{Description de la navigation}}\\
            \vspace{0.1cm}\\
            {\large \textbf{Départ de l'Herbaudière}}\\
            \begin{itemize}
                \itemsep0em 
                \item Départ de l'Herbaudière à 7h30. 
                \item La navigation peut se faire en équipage de nuit pour gagner du temps au départ et permettre de commencer le rangement et le ménage des bateaux pendant la navigation
                \item Le chenal du port de l'Herbaudière est étroit et nécessite de respecter un allignement précis ce qui peut compliquer l'accès à marée basse. En conséquence les bateaux CF encadreront les deux flotilles pour s'assurer que l'alligment soit respecté. Le départ se fait à PM-1 ce qui limite toutefois le risque. 
                \item Respecter l'alignement à 7.5° en sortant du chenal d'accès pour éviter la Potée au Beurre et les différentes sondes non balisées à l'Est du chenal jusqu'à la cardinal nord Basse du Martroger.
                \item Hissage des voiles lorsqu'on est dégagé du chenal et qu'on a de l'eau à courir sur décision du CF. 
            \end{itemize}
            {\large \textbf{Déroulé de la navigation}}\\
            \begin{itemize}
                \itemsep0em 
                \item Route vers la cardinal Nord banc de la Blanche. 
                \item Une fois la carindal enroulée cap vers le 75° vers Pornic pendant environ 7Mn
                \item Attention à bien laisser la cardinale Ouest Notre Dame à l'Est. 
                \item Attention à laisser la marque spéciale devant Pornic à l'Ouest. 
            \end{itemize}
        
        {\large \textbf{Arrivée à Pornic}}\\
            \begin{itemize}
                \itemsep0em 
                \item Affalage des voiles devant le chenal de Pornic. Attention les abords du chenal ne sont pas francs, enrouler les balises babords sans trop s'approcher de celles tribords.
                \item Objectif flotte amarrée à 10h30, rangement et ménage des bateaux.
            \end{itemize}
\end{minipage}%
%%%%%%%%%%%%%%%%
% Page de droite
%%%%%%%%%%%%%%%%
\newgeometry{top=2cm, bottom=2cm, left=2cm, right=2cm}
\begin{minipage}[t][0.9\textheight][t]{0.60\textwidth}
    \color{StrongBlue}    
    \noindent
    \begin{minipage}[t][0.3\textheight][t]{\textwidth}
        \begin{minipage}[t][0.3\textheight][t]{0.65\textwidth}
            \vspace{-8.5em}
            {\huge \textbf{Détails de la navigation}}\\
            Le QR-code situé à droite est cliquable, il permet d'explorer le planning de navigation sur le site datashom. 
            La carte ci-dessous visualise le trajet prévu ainsi que les abris de replis et dangers identifiés. 
        \end{minipage}%
        \hfill
        \begin{minipage}[t][0.3\textheight][t]{0.25\textwidth}       
        \href{https://data.shom.fr/datacontext/fr/1780679297167-806330}{ 
            \includegraphics[width=0.8\textwidth]{qr-code.png}
        }
        \end{minipage}%
    \end{minipage}
    \begin{minipage}[t][0.7\textheight][t]{\textwidth}
        % \includegraphics[width=0.97\textwidth]{carte.png}
        \includegraphics[height=0.7\textheight]{carte.png}
    \end{minipage}
\end{minipage}
\begin{minipage}[t][0.9\textheight][t]{0.40\textwidth}
    \color{StrongBlue} 
    \centering 
    \includegraphics[height=\textheight, width=\textwidth, keepaspectratio]{illustration_verticale.jpg}
\end{minipage}%
\restoregeometry



